\chapter{Continued Interactive GSLTs}
\label{ch:cigslt}
%===============================================================================

Interactivity names a site. The constructions of Part~III need more: they need to be
able to \emph{intervene} at the site --- to interpose a token, to record an event, to
attach a guard. Three further conditions suffice, and this section isolates them.

\section{Section-equipped}

The first condition concerns the relationship between a term and its congruence class.
Let $T$ be the free syntax of the interacting sort and $\widehat{T} = T/{\equiv}$ its
quotient by structural congruence. Two maps relate them:

\begin{itemize}[leftmargin=2em]
  \item the \emph{quotient} $T \twoheadrightarrow \widehat{T}$, respecting $\equiv$,
    invertible up to congruence. This is the structure a name constructor exposes when
    channels are taken up to $\equiv$ --- in rho, $@P$ with $@P = @Q \iff P \equiv Q$.
  \item a \emph{section} $\cf : \widehat{T} \to T$, a choice of canonical representative
    per class.
\end{itemize}

\begin{definition}[Section-equipped]
\label{ob:def:section}
An interactive GSLT is \textbf{section-equipped} when it comes with a computable section
$\cf : \widehat{T} \to T$ of its $\equiv$-quotient --- equivalently, a computable
canonical-form function on the interacting sort.
\end{definition}

The condition is sharper than it looks, and the $\lambda$-calculus is the test that
shows why. The $\lambda$-calculus has no analogue of the quote constructor $@$. It
passes anyway, because what is required is not a channel constructor but a section: for
$\lambda$ the congruence is $\alpha$ and a canonical representative of each
$\alpha$-class is the de Bruijn encoding~\cite{deBruijn1972}. Rho fuses the quotient and
the section into one operator viewed two ways; $\lambda$ pulls them apart; the
requirement is identical in both.

\section{Wrappability}

The second condition asks that the contraction tolerate having its operands decorated.
Anticipating Chapter~\ref{ch:costmonad}, write $\Tw$ for a sort of \emph{wrapped} terms.

\begin{definition}[Wrappability]
\label{ob:def:wrappable}
The contraction $\compute$ is \textbf{wrappable} when it is well-sorted as an operation
on wrapped continuations,
$\compute : \Surf \times \Surf \times \Tw \times \Tw \to \Tw \times \Tw$ ---
equivalently, when $\compute$ is definable on decorated continuations, mapping decorated
inputs to decorated outputs.
\end{definition}

The condition is mild for substitution-like contractions: substituting a wrapped term
into a position yields a wrapped term, so $\compute$ preserves the sort for free and
need not know that wrapping exists. It excludes contractions that inspect and dismantle
their operands --- an opaque contraction returning bare residuals would violate it.

\section{The definition}

\begin{definition}[Continued interactive GSLT]
\label{ob:def:cigslt}
A \textbf{continued interactive GSLT} is an interactive GSLT equipped with:
\begin{enumerate}[label=(\roman*),leftmargin=2.5em]
  \item a presentation of its dynamics in interaction-cut form $(\dagger)$, naming
    $K, K_p, K_e, K'$ and the partial contraction $\compute$;
  \item a section (Definition~\ref{ob:def:section}); and
  \item wrappability (Definition~\ref{ob:def:wrappable}).
\end{enumerate}
We write $\catciGSLT$ for the resulting category, whose morphisms are theory maps that are
bisimulation-preserving on the interacting sort and \emph{quote-faithful}: injective on
the reachable signature keys.
\end{definition}

The adjective names a strict strengthening. Clauses (i)--(iii) are data and conditions
added on top of an interactive GSLT, not a restriction of it; forgetting them recovers
the underlying interactive GSLT.

\begin{definition}[The forgetful functor]
$U : \catciGSLT \to \catiGSLT$ sends a continued interactive GSLT to its underlying interactive
GSLT, discarding the cut presentation, the section, and the wrappability witness, and
acts as the identity on the underlying theory map of a morphism.
\end{definition}

\begin{proposition}[The distinction is proper]
\label{ob:prop:proper}
$U$ is faithful but neither full nor essentially surjective. Hence $\catciGSLT$ is a
genuine, proper subcategory of $\catiGSLT$: strictly more structure on objects, strictly
fewer morphisms between them.
\end{proposition}

\begin{proof}[Discussion]
\emph{Faithful, not full.} $U$ forgets no information about a morphism's action, so it
is faithful. It is not full because a $\catciGSLT$-morphism must additionally be
quote-faithful: an $\catiGSLT$-morphism that is bisimulation-preserving yet collapses two
distinct signature keys is a map between the underlying objects with no lift.

\emph{Not essentially surjective.} Clause (i) requires the dynamics to factor as contact
then contraction, which a theory with non-local or unfactorable rewrites need not;
clause (ii) requires a computable section, which a theory with undecidable structural
congruence lacks; clause (iii) requires the contraction to preserve wrapping. An
interactive GSLT failing any of these has no preimage.
\end{proof}

\begin{remark}[$\lambda$ exhibits the distinction]
The $\lambda$-calculus is an object of $\catiGSLT$ as such: terms, $\alpha$, $\beta$. It
becomes an object of $\catciGSLT$ only once one \emph{chooses} the added structure --- the
de Bruijn section for (ii), the degenerate-$K_e$ presentation of $\beta$ for (i),
substitution-on-thunks for (iii). The same underlying calculus sits at both levels,
witnessing that ``continued'' is added structure and not a property read off the
interactive GSLT alone.
\end{remark}

\section{How the examples are examples}
\label{ob:subsec:table}

\begin{remark}[The migration spectrum]
\label{ob:rem:migration}
The instances are usefully ordered by how much information the contraction moves from
the environment continuation into the program continuation. Four modes occur: no
migration (CCS, interaction categories), migration by binding (rho, $\pi$, $\lambda$),
migration by spatial restructuring (ambients), and migration as interface composition.
\end{remark}

Table~\ref{ob:tab:cut} records the interaction-cut data of each running example, and
Table~\ref{ob:tab:verdict} records whether that data satisfies the successive conditions.
The final column of Table~\ref{ob:tab:verdict} is the discrimination the two strengthenings
were introduced to make.

\begin{table}[t]
\footnotesize
\centering
\begin{tabular}{@{}lllllp{34mm}@{}}
\toprule
\textbf{Instance} & \textbf{$K$} & \textbf{$E$ on $K$} & \textbf{$I$, $J$} &
\textbf{$C_p$ / $C_e$} & \textbf{$\compute$ (migration)} \\
\midrule
Calculator & --- & --- & --- & --- & no interaction rule \\[2pt]
CCS & ${\mid}$ & AC & $a$, $\bar a$ & $P$ / $Q$ &
  pure release $P \mid Q$ (none) \\[2pt]
Interaction cats. & $\circ$ & --- & shared interface & proc. / proc. &
  interface composition (none) \\[2pt]
rho, $\pi$ (sync) & ${\mid}$ & AC & $x$, $x$ & $\lambda y.P$ / $[Q_1,Q_2]$ &
  pseudo-app.\ $P\{Q_1/y\} \mid Q_2$ (binding) \\[2pt]
rho, $\pi$ (async) & ${\mid}$ & AC & $x$, $x$ & $\lambda y.P$ / $[Q,-]$ &
  pseudo-app.\ $P\{Q/y\}$ (binding) \\[2pt]
$\lambda$ & $\mathrm{App}$ & free & $x$, $-$ & $\lambda x.M$ / $N$ &
  application $M\{N/x\}$ (binding) \\[2pt]
Ambients & ${\mid}$ & AC & $n$, $n$ & $P$ / $Q$ &
  boundary dissolve/move (spatial) \\[2pt]
\midrule
Turing & \emph{none} & --- & --- & --- & control/tape heterogeneous \\[2pt]
Moore & \emph{none} & --- & --- & --- & control/stream heterogeneous \\[2pt]
Mealy & \emph{none} & --- & --- & --- & control/stream heterogeneous \\
\bottomrule
\end{tabular}
\caption{Interaction-cut data. A surface entry of $-$ denotes a nominally degenerate
surface carried structurally by the rigidity of $K$; an environment-continuation entry
of $-$ denotes a degenerate continuation, which is what distinguishes asynchronous from
synchronous output.}
\label{ob:tab:cut}
\end{table}

\begin{table}[t]
\footnotesize
\centering
\begin{tabular}{@{}lccccl@{}}
\toprule
\textbf{Instance} & \textbf{GSLT} & \textbf{has $K$-rule} & \textbf{section $\cf$} &
\textbf{wrappable} & \textbf{Verdict} \\
\midrule
Calculator      & \checkmark & --- & normal form & --- & GSLT only \\
CCS             & \checkmark & \checkmark & normal form for AC $\mid$ & \checkmark & continued interactive \\
Interaction cats.& \checkmark & \checkmark & interface normal form & \checkmark & continued interactive \\
rho, $\pi$ (sync) & \checkmark & \checkmark & normal form for AC $\mid$ & \checkmark & continued interactive \\
rho, $\pi$ (async)& \checkmark & \checkmark & normal form for AC $\mid$ & \checkmark & continued interactive \\
$\lambda$        & \checkmark & \checkmark & de Bruijn & \checkmark & continued interactive$^\ast$ \\
Ambients         & \checkmark & \checkmark & normal form for AC $\mid$ & \checkmark & continued interactive \\
\midrule
Turing           & \checkmark & --- & --- & --- & GSLT only \\
Moore            & \checkmark & --- & --- & --- & GSLT only \\
Mealy            & \checkmark & --- & --- & --- & GSLT only \\
\bottomrule
\end{tabular}
\caption{Verdicts. $^\ast$ $\lambda$ qualifies only once the added structure is chosen;
see Remark above. The three machines are perfectly good GSLTs and fail at the first
strengthening, not the second.}
\label{ob:tab:verdict}
\end{table}

Two readings of the tables are worth making explicit.

First, the calculator row and the three machine rows fail for \emph{different} reasons,
and the tables record the difference. The calculator has no interaction rule because it
has no notion of a program meeting an environment at all --- it is a language of closed
expressions. The machines have a notion of the control meeting something, but the
something is of the wrong sort. Both end at ``GSLT only,'' and a reader who wanted to
know which of the two situations obtains would have to look at the $K$ column.

Second, $\lambda$ is the only row whose verdict carries a caveat, and the caveat is the
content of Proposition~\ref{ob:prop:proper}. Nothing about the $\lambda$-calculus as
ordinarily presented determines a section; de Bruijn is a choice, and a different choice
--- a different canonical form for $\alpha$-classes --- would give a different object of
$\catciGSLT$ over the same object of $\catiGSLT$. This is what it means for $U$ to fail to be
essentially surjective on the nose while being surjective on the underlying calculi one
cares about.

